\section{Literature Review}\label{sec:literature}

Spatial competition provides a foundation for understanding how firms'
positions influence demand and the profitability of their decisions.
In \citet{hotelling_stability_1929}, consumers choose between firms by
weighing prices against transportation costs, making location an important
determinant of competitive advantage.
A firm's decision to reposition consequently depends on the customers it can
attract and the responses of its competitors.
These interactions become more complex when firms operate across several
markets. \citet{bulow_geanakoplos_klemperer_multimarket_1985} establish how
strategic substitutes and complements connect competitive decisions, while
\citet{bimpikis_ehsani_ilkilic_cournot_2019} examine how competition and
production linkages transmit shocks through networked markets.
Together, these contributions motivate studying firm positioning through the
allocation of resources across related activities and the competitive
responses those allocations generate.

The multi-contest framework of \citet{xu_zenou_zhou_conflict_2022} provides
the closest theoretical foundation for our analysis.
In their model, players allocate costly effort across contests, receiving
shares of prizes that depend on their allocations relative to those of competitors.
Their equilibrium characterization and comparative statics establish how
competitive interactions transmit changes in incentives through a network of contests.
We adapt this proportional-share structure to firms allocating operating
capacity across activities that serve overlapping customer demands.
Capacity determines a firm's ability to serve each demand opportunity, while
rivals compete for the same available contribution margins.
The resulting allocation problem admits a discrete-choice interpretation when
capacity represents varieties or outlets
\citep{mcfadden_conditional_logit_1974,train_discrete_choice_2009}.
Building on the equilibrium and sensitivity methods of
\citet{xu_zenou_zhou_conflict_2022}, we establish uniqueness and characterize
firms' responses to changes in market opportunities and costs under the
conditions of our capacity model.

Competitive repositioning also depends on the costs of changing existing
investments. \citet{sweeting_positioning_2013} studies forward-looking
product positioning in differentiated markets, demonstrating the importance
of adjustment costs for firms' decisions.
We introduce persistence through quadratic costs of changing inherited
capacity, with firms making sequential Nash choices as market opportunities
and competitors' allocations evolve.
The resulting operating path determines the cash flows investors anticipate
following an exogenous free cash flow shock. \citet{berk_green_naik_growth_options_1999}
relate investment opportunities to security returns through changes in assets
and growth options.
We follow this connection between operating opportunities and financial
valuation, discounting the cash flows generated along the adjustment path.
Firm valuations can consequently respond when new information arrives, while
firms' operating allocations adjust over time.

Corporate disclosures and financial news provide information about firms'
activities and their competitive environments. \citet{hoberg_phillips_2016}
demonstrate how corporate text can reveal product-market relationships and
differences in competitive positioning, while
\citet{gentzkow_text_as_data_2019} establish methodological foundations for
using text as economic data.
The information available to investors also reflects the incentives governing
its production. \citet{gentzkow_media_bias_2006} study how reputation and
reporting incentives influence news content, providing an economic motivation
for considering differences in the coverage of firms' activities.
We incorporate these considerations through an observation model in which a
firm's news distribution depends on both its operating capacity and the
relative intensity of reporting across activities.

Representing corporate information as a distribution preserves variation
among the activities and perspectives reflected in individual articles.
\citet{kusner_word_embeddings_2015} develop transport-based comparisons of
word embeddings, while \citet{gawronsky_continuous_2026} use Wasserstein
geometry to study covariance envelopes from language-model representations.
Optimal transport compares distributions through the minimum cost of matching
their probability mass, retaining information about their composition and
dispersion that may be lost when articles are represented by a single average vector.
This provides a natural basis for comparing firms' positions in information
space and measuring changes in those positions over time.

We extend this literature by deriving an economic mechanism for changes in
firms' economic footprints.
Costly capacity adjustment and competitive responses determine how firms
react to changing opportunities, while financial valuation reflects the
cash-flow consequences of those decisions.
In a restricted two-opportunity economy, initial exposure determines the
ordering of firms' return responses and, under the specified reporting
conditions, the direction of movement toward other firms' initial news distributions.
The resulting nonnegative expected local association provides the theoretical
basis for the estimator developed in the following section.
